\chapter{Introducción}
\label{chap:introduccion}

\lettrine{E}{l} abandono temprano de la educación tiene importantes implicaciones sociales en términos de calidad de vida, salud, marginalidad social, delincuencia o empleabilidad. Dado el importante coste socioeconómico que este fenómeno conlleva a largo plazo, todas las administraciones públicas y organismos gubernamentales tratan de reducir en lo posible la tasa de abandono temprano y fracaso escolar mediante políticas de intervención y apoyo proactivo \cite{OCDE_Educacion}.

\section{Contexto}

La administración pública, habitualmente a nivel de comunidad autónoma, es la entidad encargada de la gestión de la educación y debe enfrentarse al complejo desafío de realizar un reparto equitativo de los escasos recursos materiales y personales entre su extensa red de centros educativos. 

La casuística de cada institución es única, y algunos centros necesitan un volumen de apoyo significativamente mayor que otros debido a múltiples factores sistémicos y estructurales. Estar ubicados en zonas económicamente deprimidas, contar con un alto porcentaje de alumnado proveniente de clases desfavorecidas o enfrentarse a severas barreras de brecha digital son elementos que lastran el rendimiento global de un centro. Esto se alinea con la realidad del indicador AROPE en España, que evidencia cómo un porcentaje significativo de las familias conviven en situaciones de vulnerabilidad o pobreza relativa \cite{INE_Pobreza_PDF, AyudaEnAccion_Pobreza}. Abordar estas carencias desde el prisma de la "pobreza multidimensional" \cite{BBVA_Pobreza_Multidimensional, Observatorio_Pobreza} es fundamental para que la administración actúe como un equilibrador social efectivo.

Sin embargo, históricamente la información ha estado dispersa en bases de datos heterogéneas a nivel gubernamental. Para hacer un reparto verdaderamente objetivo y eficiente, la administración no puede depender de intuiciones o reportes fragmentados; debe disponer de una "fotografía" exacta del estado de cada centro, construida sobre medidas estandarizadas, comparables y objetivas a nivel autonómico.

\section{Motivación}

Ante este escenario de gran volumen y dispersión de datos, surge la necesidad imperativa de aplicar técnicas de Inteligencia de Negocios (\textit{Business Intelligence} o BI) a la gestión pública. Dado el inmenso número de centros educativos que suelen estar bajo la tutela de las administraciones, la toma de decisiones manual y aislada resulta inabarcable.

La motivación fundamental de este Trabajo de Fin de Grado es el diseño y desarrollo de un ecosistema analítico que culmine en la creación de un cuadro de mando. Dicho cuadro de mando estará compuesto por una serie de Indicadores Clave de Desempeño (KPIs, por sus siglas en inglés) que sinteticen y crucen de forma inteligente el rendimiento académico con el contexto sociodemográfico de los centros a escala global. 

El propósito final es que la tecnología se ponga al servicio de los gestores públicos, sirviendo como una herramienta centralizada y accionable para guiar la difícil tarea de asignación de recursos, distribución de presupuestos y dotación de personal de apoyo entre los distintos centros educativos de forma totalmente empírica y justificada.

\section{Objetivos}
\label{sec:objetivos}

El objetivo principal de este proyecto es crear un almacén de datos (Data Warehouse) que extraiga su información a partir de diversas fuentes institucionales, y el desarrollo de un cuadro de mando que muestre distintos KPIs de interés estratégico para los gestores públicos.

Para alcanzar este propósito, garantizando además que todo el sistema se diseñe utilizando exclusivamente herramientas libres de licencias de pago, el trabajo se ha desglosado en los siguientes objetivos específicos:

\begin{enumerate}
    \item \textbf{Simulación de fuentes de datos:} Crear tres bases de datos de origen que representen el entorno autonómico: una con expedientes académicos de alumnos, una con adaptaciones curriculares y una tercera con los resultados de encuestas socioeconómicas familiares.
    \item \textbf{Diseño del almacén de datos:} Modelar un repositorio centralizado bajo un enfoque dimensional (Datamart) que aglutine e integre la información de las tres fuentes mencionadas para ser el origen analítico de los KPIs.
    \item \textbf{Diseño de procesos ETL:} Diseñar e implementar los flujos de Extracción, Transformación y Carga (ETL) que pueblen los datos de las fuentes de origen hacia el almacén de datos.
    \item \textbf{Desarrollo de la lógica de negocio (Backend):} Desplegar una API robusta y segura que sirva como capa de servicio estandarizada entre el almacén de datos y la interfaz gráfica.
    \item \textbf{Creación de la aplicación web (Dashboard):} Implementar el cuadro de mando interactivo destinado a los gestores, diseñado para mostrar y explorar visualmente los KPIs desglosados por centro, localidad y período de tiempo.
\end{enumerate}

De este modo, el proyecto abarca de manera secuencial e iterativa la totalidad del ciclo de vida de un proyecto de Inteligencia de Negocios aplicado a la administración pública.
